%% ===== Back matter: declarations + bibliography =====
\backmatter

\bmhead{Acknowledgments}

\section*{Declarations}
% Funding; competing interests; author contributions.
% (Data availability + Code availability live at the END of Methods, per Nature style.)

%% ---- Supplementary information: SI figures grouped here, references placed AFTER them.
%% \clearpage flushes pending floats so the short bibliography can't interleave with the SI
%% figures -- that interleaving was the figure/reference clash. ----
\clearpage
\bmhead{Supplementary information}
% Supplementary Figures S1-S6 (see outline). S1 = classical-ML baselines.

% Supplementary figures use S-numbering. The TorchCell overview schematic is parked
% here for now (was Fig 1 in the main text); offset the counter so it renders as
% Suppl. Fig. S7 and the planned S1-S6 keep their numbers. Renumber when the SI is
% built out. Same true-to-size, left-justified, caption-matched-to-width layout.
\setcounter{figure}{6}
\renewcommand{\thefigure}{S\arabic{figure}}

% Figures are placed true-to-size with \tcfig (NO LaTeX scaling): each draw.io source is
% drawn inside the 180x170 mm print-box boundary template, so the --crop export equals the
% Nature print box and renders 1:1. If a figure exports wider than 180 mm, fix it in draw.io
% (pull content inside the box) -- do not scale. Order/numbering is provisional.
\begin{figure}[htbp]
\tcfig{figures/TorchCell-Supervised-Learning-and-Teacher-Forcing-Generic-Phenotypes.pdf}{TorchCell: ontology-unified knowledge graph and the
perturbation-operator cell representation. Panels (a)--(d).}{fig:overview}
\end{figure}

\begin{figure*}[t]
\tcfig{figures/ontology_pydantic_hourglass_data_model.pdf}{Ontology data model (hourglass). Each experiment
funnels genotype, phenotype, environment, and experiment-specific metadata through a
shared Data Merge Layer that distinguishes records from the same experiment; the layer
can be ignored when merging into larger datasets.}{fig:data-model}
\end{figure*}

\begin{figure*}[t]
\tcfig{figures/neo4j_cell-conversion-deduplication-aggregation.pdf}{Knowledge-graph normalization in three
stages. Conversion maps phenotype experiment types (GeneEssentiality, SyntheticLethality)
to a common Fitness type; deduplication collapses records equivalent other than their
phenotype-label values; aggregation combines equivalent-genotype records with different
phenotype labels (Fitness, GeneInteraction) into a single HeteroData instance.}{fig:conversion-dedup-agg}
\end{figure*}

\begin{figure*}[t]
\tcfig{figures/s288c_selecting_gene_sequence.pdf}{DNA sequence selection for the sequence model.
The coding sequence is taken from the outermost start to stop codon, ignoring five-prime
UTR introns unless internal to the gene. Cases: internal intron (YFL039C), 5' UTR intron
(YBL072C), and a 1\,bp CDS lacking a start codon (YIL111W).}{fig:dna-selection}
\end{figure*}

\begin{figure*}[t]
\tcfig{figures/fungal_up_down_transformer_upstream_1003bp_pad_for_undersized.pdf}{DNA tokenization. A 1003\,bp window is
encoded as [CLS], species proxy, 6-mer sequence, [SEP]; undersized sequences are padded
with [PAD] to a fixed length before vectorization.}{fig:dna-tokenization}
\end{figure*}

\begin{figure*}[t]
% TODO (draw.io source uncharacterized-genes-profile-enricment.drawio.png):
%   add (f) total quantification of triple interactions across the remaining
%     (uncharacterized) genome -- how many triple interactions exist in yeast / coverage;
%   add (g) the new GO annotations produced by the guilt-by-association transfer.
%   Pull all content inside the 180x170 mm boundary box (currently overflows to ~214 mm wide).
\tcfig{figures/uncharacterized-genes-profile-enricment.pdf}{Guilt-by-association annotation of
uncharacterized genes. (a)--(e) Triple-interaction profiles $p$ are correlated (Pearson, PCC)
across all gene pairs; pairs with PCC $>0.6$ are linked, and GO enrichment over the
resulting network transfers annotations ($\mathrm{GO}_{g_i}=\mathrm{GO}_{g_j}$).
(f) Total quantification of triple interactions across the remaining yeast genome
(coverage of uncharacterized genes). (g) The resulting new annotations transferred to
previously uncharacterized genes.}{fig:profile-enrichment}
\end{figure*}

%% ---- References: placed AFTER the SI part. \clearpage first so every SI figure float is
%% output before the bibliography begins (prevents the float/reference clash). NOTE: for
%% Springer Nature submission, paste the .bbl here instead of \bibliography{} (see `make flat`). ----
\clearpage
\nocite{*}% FILLER: force all seeded refs to render in the preview; remove for submission
\bibliography{references}
